\section{Every Entity Gets a Name}
\label{sec:ch04-every-entity-gets-a-name}

\index{AddGlobalObjectIdentification@\code{AddGlobalObjectIdentification}}%
Two calls go on the end of \code{Program.cs}, and the second one belongs to
section~\ref{sec:ch04-a-mutation-that-can-fail}. They are the only lines in
the file that chapter~\ref{ch:data-without-n-plus-1} did not already print;
the line above them gave up its semicolon to make room, and nothing else
moved.

\begin{minted}[highlightlines={32,33}]{csharp}
using ChilliCream.Nitro.App;
using HotChocolate.AspNetCore;
using Microsoft.EntityFrameworkCore;
using Sessions;

var builder = WebApplication.CreateBuilder(args);

// EF Core reports every command through the logging pipeline as well,
// at Information and with a duration attached. StatementLog below is
// this service's instrument, so silence the other one rather than have
// every statement printed twice.
builder.Logging.AddFilter(
    "Microsoft.EntityFrameworkCore.Database.Command",
    LogLevel.Warning);

builder.Services.AddDbContextFactory<ConferenceContext>(options =>
{
    options.UseSqlite("Data Source=conference.db");

    if (builder.Environment.IsDevelopment())
    {
        options.AddInterceptors(new StatementLog());
    }
});

builder
    .AddGraphQL()
    .AddSessionsTypes()
    .RegisterDbContextFactory<ConferenceContext>()
    .AddQueryContext()
    .AddPagingArguments()
    .AddGlobalObjectIdentification()
    .AddMutationConventions();

var app = builder.Build();

using (var scope = app.Services.CreateScope())
{
    var factory = scope.ServiceProvider
        .GetRequiredService<IDbContextFactory<ConferenceContext>>();
    using var db = factory.CreateDbContext();
    db.Database.EnsureCreated();
}

app.MapGraphQL()
    .WithOptions(options =>
    {
        // Serve the Nitro build that shipped in the package. The
        // default, ServeMode.Latest, downloads the current one
        // from ChilliCream's CDN instead.
        options.Tool.ServeMode = ServeMode.Embedded;
    });

app.RunWithGraphQLCommands(args);
\end{minted}

\index{NodeResolver@\code{[NodeResolver]}}%
Then one method per type that wants to be reachable. \code{SessionType.cs}
gains a second resolver beside the speaker one, and the new method is the same
query \code{sessionById} already runs. It also gains two \code{using}
directives, because until now this file had none:

\begin{minted}[highlightlines={1,2,18-30}]{csharp}
using GreenDonut.Data;
using Microsoft.EntityFrameworkCore;

namespace Sessions;

[ObjectType<Session>]
public static partial class SessionType
{
    /// <summary>
    /// The person giving this session.
    /// </summary>
    public static async Task<Speaker?> GetSpeakerAsync(
        [Parent(nameof(Session.SpeakerId))] Session session,
        ISpeakerByIdDataLoader speakerById,
        CancellationToken ct)
        => await speakerById.LoadAsync(session.SpeakerId, ct);

    /// <summary>
    /// Loads one session by the identifier inside a global node id.
    /// </summary>
    [NodeResolver]
    public static async Task<Session?> GetSessionAsync(
        int id,
        ConferenceContext db,
        QueryContext<Session> query,
        CancellationToken ct)
        => await db.Sessions
            .Where(s => s.Id == id)
            .With(query)
            .FirstOrDefaultAsync(ct);
}
\end{minted}

\index{SpeakerType@\code{SpeakerType}!exists only to be a node}%
\code{Speaker} has never needed a type class, and now it does. The whole of
\code{SpeakerType.cs} is the node resolver, and it borrows
chapter~\ref{ch:data-without-n-plus-1}'s DataLoader rather than writing a
query of its own:

\begin{minted}{csharp}
namespace Sessions;

[ObjectType<Speaker>]
public static partial class SpeakerType
{
    /// <summary>
    /// Loads one speaker by the identifier inside a global node id.
    /// </summary>
    [NodeResolver]
    public static async Task<Speaker?> GetSpeakerAsync(
        int id,
        ISpeakerByIdDataLoader speakerById,
        CancellationToken ct)
        => await speakerById.LoadAsync(id, ct);
}
\end{minted}

\index{node resolver!takes the decoded key}%
Both take a plain \code{int}. The decoding happens before the method is
entered, so what arrives is the key that came out of the string, and the
resolver never sees an id it has to parse. That is why neither method carries
an \code{[ID]} attribute: attaching the node resolver is what declares the
type a node, and being a node is what rewrites the id field.

\index{schema!what a node resolver generates}%
Nothing else changed, and the exported schema gained all of this. As in
chapter~\ref{ch:data-without-n-plus-1}, the descriptions and the \code{@cost}
directives the export writes are dropped here and the unchanged fields of the
two types are elided, because the shape is what the section is about:

\begin{graphqlsdl}
type Query {
  node(id: ID!): Node
  nodes(ids: [ID!]!): [Node]!
  ...
}

interface Node {
  id: ID!
}

type Session implements Node {
  id: ID!
  ...
}

type Speaker implements Node {
  id: ID!
  ...
}
\end{graphqlsdl}

\index{Node@\code{Node}!the book's first interface}%
That is the book's first interface, and nothing in this chapter asked for one.
An interface is a promise that every type implementing it carries those
fields, and the only reason \code{node} can exist is that \code{Node} makes
the promise: a field returning an interface is a field whose concrete type the
client discovers at run time, which is why every query against \code{node}
carries either \code{\_\_typename} or an inline fragment.
\index{inline fragment!definition}%
An inline fragment is the \code{... on Session \{ title \}} form, and it
selects its fields only when the object turns out to be that type; the
subsection after next sends one.

\index{nodes@\code{nodes}!is Hot Chocolate's own}%
\code{nodes} is Hot Chocolate's, not the specification's. The specification
describes plural identifying root fields in general and names none of them, so
the field is a convention with a good reason and not a requirement.

\index{node id!what it decodes to}%
Send the paged query and read the identifiers instead of the names.

\begin{minted}{text}
POST /graphql HTTP/1.1
Host: localhost:5001
Content-Type: application/json

{"query":"{ sessions(first: 10) { nodes { id speaker { id } } } }"}
\end{minted}

\begin{minted}[fontsize=\footnotesize]{json}
{"data":{"sessions":{"nodes":[
  {"id":"U2Vzc2lvbjox","speaker":{"id":"U3BlYWtlcjox"}},
  {"id":"U2Vzc2lvbjoy","speaker":{"id":"U3BlYWtlcjox"}},
  {"id":"U2Vzc2lvbjoz","speaker":{"id":"U3BlYWtlcjoy"}},
  {"id":"U2Vzc2lvbjo0","speaker":{"id":"U3BlYWtlcjoz"}}]}}}
\end{minted}

Decode the first two:

\begin{minted}{text}
U2Vzc2lvbjox  ->  Session:1
U3BlYWtlcjox  ->  Speaker:1
\end{minted}

\index{node id!format}%
The type name, a colon, the key, and nothing else. No length prefix, no
version byte, no signature. Both objects are still row one of their table and
the string now says which table, which is the entire fix and the reason
section~\ref{sec:ch04-the-schema-is-the-version}'s collision cannot recur.

\index{statement count!node field}%
One \code{node} costs one statement, and the projection still applies, because
that resolver took a \code{QueryContext<Session>} like every other query in
chapter~\ref{ch:data-without-n-plus-1}. Ask \code{node} for a title and SQLite
is asked for two columns.

\begin{minted}{text}
-- statement 1
SELECT "s"."Id", "s"."Title"
FROM "Sessions" AS "s"
WHERE "s"."Id" = @id
LIMIT 1
\end{minted}

\index{nodes@\code{nodes}!batches nothing}%
\code{nodes} costs one statement per id. Two ids of two types cost two
statements, and so do two ids that are both sessions, which one query could
obviously have fetched together. The field takes a list and resolves it one
element at a time, which is the shape
chapter~\ref{ch:data-without-n-plus-1} named and did not fix here. A node
resolver that goes through a DataLoader gets the batching anyway, which is
what \code{SpeakerType} above quietly has and \code{SessionType} does not.

\subsection{Three Ways to Get an Id Wrong}

\index{node id!malformed}%
A string that is not a node id at all answers 200, resolves the field to null,
and reports an error carrying back what you sent:

\begin{minted}[fontsize=\footnotesize]{json}
{"errors":[{"message":"The node ID string has an invalid format.",
"path":["node"],"extensions":{"originalValue":"not-an-id"}}],
"data":{"node":null}}
\end{minted}

\index{node id!rejected at coercion}%
The second is a bad id that never reaches execution at all.
Section~\ref{sec:ch04-a-mutation-that-can-fail}'s mutation takes its session
as an \code{ID!} inside an input type, and the string \code{2} is a perfectly
legal \code{ID} that is not a node id. It answers 400, with no \code{data} key
anywhere in the response:

\begin{minted}[fontsize=\footnotesize]{json}
{"errors":[{"message":"The node ID string has an invalid format.",
"extensions":{"originalValue":"2"}}]}
\end{minted}

Same message, same complaint, two different stages, two different statuses.
An argument on a root field is coerced as the field executes; a variable
feeding an input object is coerced before the operation runs at all. Worth
knowing before you write the client code that has to tell a user what went
wrong.

\index{node id!wrong type resolves silently}%
The third case is the one to actually worry about, because nothing fails.
Feed \code{node} a well-formed speaker id under a fragment expecting a
session. The operation:

\begin{minted}{graphql}
query N($id: ID!) {
  node(id: $id) {
    __typename
    ... on Session { title }
  }
}
\end{minted}

\index{variables!definition}%
sent with \code{\{"id": "U3BlYWtlcjox"\}} as its variables, the JSON object a
client sends beside an operation to fill in the names that start with a
dollar sign. That id is speaker~1:

\begin{minted}{text}
{"data":{"node":{"__typename":"Speaker"}}}
\end{minted}

Status 200, no error key, and an object with none of the fields that were
asked for. The fragment did not match, so it contributed nothing, and that is
correct behavior rather than a defect: \code{node} promised a \code{Node} and
delivered one. It is also exactly what a mixed-up id looks like from the
client, which is why \code{\_\_typename} belongs in the query and not only in
the debugging session that follows.

\subsection{Retiring the Field This Replaces}

\index{deprecation!as the only removal mechanism}%
\code{sessionById} is now redundant, and removing it is not available. Deleting
a field from a schema breaks every client still asking for it, on the
deployment, with no warning and no migration window. What the specification
gives instead is a marker, and \code{Query.cs} is where it goes.

\begin{minted}[highlightlines={23-25}]{csharp}
using GreenDonut.Data;
using Microsoft.EntityFrameworkCore;

namespace Sessions;

[QueryType]
public static class Query
{
    /// <summary>
    /// Sessions on the schedule, earliest first.
    /// </summary>
    [UsePaging(DefaultPageSize = 2, MaxPageSize = 10)]
    public static IQueryable<Session> GetSessions(
        ConferenceContext db,
        QueryContext<Session> query)
        => db.Sessions
            .OrderBy(s => s.StartsAt)
            .With(query);

    /// <summary>
    /// One session by its identifier, or null if no session has it.
    /// </summary>
    [GraphQLDeprecated(
        "Ask for `node(id:)` instead. This field takes the database "
        + "key, which is only unique among sessions.")]
    public static async Task<Session?> GetSessionByIdAsync(
        int id,
        ConferenceContext db,
        QueryContext<Session> query,
        CancellationToken ct)
        => await db.Sessions
            .Where(s => s.Id == id)
            .With(query)
            .FirstOrDefaultAsync(ct);
}
\end{minted}

\index{deprecated@\code{"@deprecated}!is advisory}%
The field keeps working. Ask for it and it answers 200 with the session, and
the specification is explicit that this is the
point~\autocite{graphqlspec2025}:

\begin{quote}
It is still legal to include these fields in a selection set (to ensure
existing clients are not broken by the change), but the fields should be
appropriately treated in documentation and tooling.
\end{quote}

\index{introspection!definition}%
\index{introspection!hides deprecated fields by default}%
Tooling is where the change lands. Every GraphQL schema will describe itself
to a client that asks, through the \code{\_\_schema} and \code{\_\_type}
fields, and that is what introspection means. Introspect the \code{Query}
type and \code{sessionById} is not in the list; ask again with
\code{fields(includeDeprecated: true)} and it appears, with
\code{isDeprecated} true and the reason string attached. That is what makes an
editor stop offering the field, a schema check flag a new query against it, and
a code generator emit a warning. Nothing stops the request.

\index{deprecation!reason is the whole message}%
So the reason string is the entire user interface of the deprecation. It is
the only sentence the person removing the field will ever be shown, and
\enquote{deprecated} tells them nothing they did not already know. Name the
replacement. That is my rule rather than anybody's requirement, and it comes
from the number of reasons I have read that said \enquote{use the new one}.

\index{deprecation!not on a required argument}%
One place it will not go. The specification bars the directive from a required
argument, meaning one that is non-null with no
default~\autocite{graphqlspec2025}:

\begin{quote}
The \code{@deprecated} directive must not appear on required (non-null without
a default) arguments or input object field definitions. [\ldots] To deprecate a
required argument or input field, it must first be made optional by either
changing the type to nullable or adding a default value.
\end{quote}

\index{schema build!refuses an illegal deprecation}%
Hot Chocolate enforces it, and the failure is a good one. Put the attribute on
the \code{id} parameter of \code{sessionById}, which is an \code{Int!} with no
default, and the schema fails to build: a \code{SchemaException} naming the
argument, and a service that does not start. Not at run time on the first
request that happens to use it. The reasoning is worth carrying: a client
that omits an optional argument has taken your advice, and a client that omits
a required one has sent an invalid query, so deprecating a required argument
advises nobody of anything.

\begin{onhc14}
\code{AddGlobalObjectIdentification}, \code{[NodeResolver]} and
\code{[GraphQLDeprecated]} all exist on Hot Chocolate~14 and produce the same
\code{node} and \code{nodes} fields, the same \code{Node} interface and the
same \code{@deprecated}. The refusal to deprecate a required argument is
identical too, down to the wording of the message.

One line of this section's \code{SessionType.cs} does not compile there. The
node resolver takes a \code{QueryContext<Session>}, and 14 has no
\code{IQueryable} overload of \code{With} to apply it, which is the same gap
chapter~\ref{ch:data-without-n-plus-1}'s callout described. Drop the parameter
and the \code{.With(query)} line, and the resolver selects every column
instead of the two the client asked for. The statement count is unchanged.

The difference worth planning around is one that is not there.
\code{U2Vzc2lvbjox} is what 14 produces and it is what 16 produces, byte for
byte, for the same entity, so a client holding cached global ids keeps them
across the upgrade. That is not something you can assume of an id format in
general.

Two things do differ in what comes back, and neither is chapter~4's doing. On
14 the malformed-id error above carries an extra \code{locations} entry. And
the exported schema names a different specification URL for \code{DateTime}
and begins with a byte order mark, both of which
chapter~\ref{ch:hot-chocolate-zero} already met.
Appendix~\ref{app:hc14} carries the full mapping.
\end{onhc14}
